% Figure: Colpitts LC oscillator, drawn with plain TikZ paths and nodes (no
% circuitikz). An inductor and a capacitive divider (C1 in series with C2) form
% the resonant tank; the tap between C1 and C2 provides the feedback fraction to
% a single amplifying device drawn as a labelled gain block. The drawing is
% schematic rather than a full bias network, matching the vignette's level.
\begin{figure}[htbp]
\centering
\begin{tikzpicture}[
  line width=0.8pt,
  font=\scriptsize,
]
% ===== nodes =====
\coordinate (top)  at (0,3.4);     % top of tank (output node)
\coordinate (tap)  at (0,1.7);     % capacitive-divider tap (feedback)
\coordinate (bot)  at (0,0);       % ground rail
\coordinate (Ltop) at (2.6,3.4);
\coordinate (Lbot) at (2.6,0);

% ===== inductor branch (right leg of the tank) =====
\draw (top) -- (Ltop);
\draw (Ltop) -- (2.6,2.9);
\draw (2.6,2.9)
  .. controls (2.15,2.75) and (2.15,2.45) .. (2.6,2.3)
  .. controls (2.15,2.15) and (2.15,1.85) .. (2.6,1.7)
  .. controls (2.15,1.55) and (2.15,1.25) .. (2.6,1.1);
\draw (2.6,1.1) -- (Lbot);
\node[anchor=west] at (2.75,2.3) {$L$};
\draw (Lbot) -- (bot);

% ===== capacitor C1 (top of divider) =====
\draw (top) -- (0,2.6);
\draw[line width=1.4pt] (-0.32,2.6) -- (0.32,2.6);
\draw[line width=1.4pt] (-0.32,2.4) -- (0.32,2.4);
\draw (0,2.4) -- (tap);
\node[anchor=east] at (-0.4,2.5) {$C_1$};

% ===== capacitor C2 (bottom of divider) =====
\draw (tap) -- (0,0.9);
\draw[line width=1.4pt] (-0.32,0.9) -- (0.32,0.9);
\draw[line width=1.4pt] (-0.32,0.7) -- (0.32,0.7);
\draw (0,0.7) -- (bot);
\node[anchor=east] at (-0.4,0.8) {$C_2$};

% ===== tap node and feedback take-off =====
\filldraw[engorange] (tap) circle [radius=1.6pt];
\node[engorange, anchor=west] at (0.15,1.55) {tap};

% ===== amplifier block =====
\draw[fill=white] (4.4,1.2) rectangle (5.8,2.6);
\node at (5.1,1.9) {gain $A$};
% input from tap
\draw[engblue] (tap) -- (1.0,1.7) -- (1.0,0.55) -- (6.9,0.55) -- (6.9,1.9) -- (5.8,1.9);
\node[engblue, anchor=south, font=\tiny] at (3.9,0.57) {feedback $\beta = C_1/(C_1+C_2)$};
% output back to top of tank
\draw[engred] (4.4,1.9) -- (3.6,1.9) -- (3.6,3.9) -- (top) -- (top);
\draw[engred] (3.6,3.9) -- (top);
\node[engred, anchor=south, font=\tiny] at (1.3,3.55) {output to tank};

% ===== ground rail and symbol =====
\draw (bot) -- (Lbot);
\draw (1.3,0) -- (1.3,-0.25);
\draw (1.1,-0.25) -- (1.5,-0.25);
\draw (1.17,-0.37) -- (1.43,-0.37);
\draw (1.25,-0.49) -- (1.35,-0.49);

\end{tikzpicture}
\caption[Colpitts oscillator]{The Colpitts LC oscillator. An inductor $L$ in
parallel with the series pair $C_1$ and $C_2$ forms the resonant tank; the tap
between the two capacitors returns a fraction $\beta = C_1/(C_1+C_2)$ of the
tank voltage to the amplifier input. The loop oscillates at the tank resonance
$f_0 = 1/(2\pi\sqrt{L\,C_1 C_2/(C_1+C_2)})$ once the amplifier gain makes the
round-trip loop gain reach unity with zero net phase, the Barkhausen
condition worked in the vignette of Section~\ref{sec:vol04ch09:colpitts}.}
\label{fig:vol04ch09:colpitts}
\end{figure}
